\section*{Discussion}
With the rise of video games and generative AI technologies, humans are exposed to virtual faces daily. 
Yet, how our brain processes these faces, and their emotions has received little attention. 
Here we explored neural responses to presentations of different face-types (real, virtual), across a range of emotions (Anger, Disgust, Fear, Joy, Sadness, Surprise, Neutral control). 
A preliminary behavioural experiment confirmed that the emotional expressions of the two datasets were matched in intensity. 
Hemodynamic activity (HbT) was measured with functional near-infrared spectroscopy (fNIRS) and evaluated in two complementary ways: Activation analysis (GLM): estimated change in HbT amplitude in each channel, quantifying how strongly each Face-type \texorpdfstring{$\times$}{x} Emotion combination modulated local cortical activity; and Spectral functional connectivity: sliding-window cross-spectral density was used to compute coherence between channel pairs, revealing how neural signals became temporally synchronized (or desynchronized) across regions during stimulus processing. 
We hypothesized participants would exhibit differences in activation and connectivity to both different face-types and emotions. 

Distinct mechanisms were observed in the processing of face-types and emotions. 
Activation analyses indicated that differences for both occurred primarily in the occipital region, the brain's primary visual processing center. 
The occipital region's role extends beyond simple feature detection, as it processes complex visual patterns and configurations that differentiate face types and facial expressions at a perceptual level. 
These occipital responses likely reflect the visual system's ability to encode the physical differences in facial features and expressions before this information is forwarded to temporal and parietal cortices for higher-order processing. 
Our connectivity analyses revealed widespread differences in the frontal, parietal and temporal areas, in line with this view. 
The increased parietal-temporal-frontal coherence suggests greater network-level integration of these regions, becoming more functionally coupled when processing specific affective information (e.g., Fear $>$ Neutral) and face realism (real $>$ virtual).
Interestingly, this juxtaposition suggests that while the level of activation in these regions remains stable, the similarity in HbT activity suggests that these regions are prioritized for the processing of real faces over virtual faces, and differential processing for different emotions.

Our finding that virtual faces elicit greater left occipital activation compared to real faces aligns with evidence that face realism modulates visual processing in a complex manner. 
Our results likely reflect compensatory processing in functionally connected occipital regions that support the fusiform face area in response to perceptually challenging stimuli. 
This compensatory mechanism may manifest as the increased left occipital activation we observed for virtual faces, as these upstream visual areas work harder to provide adequate facial feature extraction to downstream face-processing networks.

Conversely, our functional connectivity analysis revealed widespread but distinct neural connectivity patterns associated with processing real versus virtual faces, suggesting that the realism of facial stimuli modulates the underlying brain network dynamics during face processing, supporting our first hypothesis.
We observed stronger connectivity when processing real faces across parietal, frontal, and central/temporal regions. 
The functional connectivity analysis revealed distinct emotional clusters that reflect underlying similarities in neural processing patterns and arousal characteristics, as seen in Figure \ref{fig:fc_emotion_summary_analysis}.
Anger and Joy also produced notably strong connectivity, clustering with Fear in our summary heatmap, suggesting these emotions share common neural network engagement patterns despite their different valence profiles.
This clustering pattern aligns with dimensional models of emotion that emphasize arousal as a key organizing principle, where high-arousal emotions like Fear, Anger, and Joy engage more extensive and coherent neural networks than their low-arousal counterparts \citep{ke_dynamic_2025}.
The emergence of these two distinct clusters, one characterized by high connectivity (Fear, Anger, Joy) and another by relatively lower connectivity (Disgust, Sadness, Surprise, Neutral), suggests that the brain's functional architecture organizes emotional face processing along arousal dimensions. 
This finding is particularly important for understanding how virtual characters and avatars might be designed to maximize/minimize neural engagement, as high-arousal emotional expressions appear to recruit broader brain networks regardless of face realism.

Our count of significantly different channel pairs across all emotions found minimal connectivity differences among left and right occipital ROI, while central/temporal and parietal regions varied strongly, as indicated by the asterisks and carets (Figure \ref{fig:fc_region_summary_analysis}).
We found only 20-40 significantly different channels in the occipital regions, compared to 100-120 in the central/temporal and parietal regions, indicating a structural difference in how these regions process emotional faces.
This fits with models of face perception, where early visual areas (e.g., occipital face area) feed into higher-level hubs like the fusiform face area. 
This suggests that while occipital areas are crucial for initial face processing, the emotional nuances of faces are primarily encoded in more distributed networks involving parietal and central/temporal regions, regardless of the emotion presented.
This lines up with the constructionist view of emotion, that emotional expressions are decoded in distributed brain networks rather than isolated regions \citep{barrett_solving_2006}. 

The interaction between face type and emotion revealed a divide among emotions. 
Joy, Sadness, and Surprise showed greater functional connectivity for real faces, while Fear, Anger, Disgust, and Neutral expressions elicited stronger connectivity for virtual faces. 
This split suggests that not all emotions are equally affected by face realism, some emotional expressions may be more prone to virtual facial representations, while others depend on the detailed information present in real faces.

The regional analysis of the face type by emotion interaction supported our previous functional connectivity findings. 
The occipital regions showed similar connectivity levels for both real and virtual faces across emotions, reinforcing the notion that early visual processing is relatively insensitive to face realism.
In contrast, the central/temporal and parietal regions exhibited greater connectivity for real faces across emotions. 
The frontal regions, however, showed the opposite pattern, with greater connectivity for virtual faces.
This increased frontal engagement for virtual faces may reflect compensatory mechanisms required to interpret/contextualize emotional expressions in less naturalistic stimuli.
These regional differences suggest that while occipital areas process both face types similarly, higher-order regions differentiate between real and virtual faces in ways that depend on the specific emotional content being processed.

This study was the first to use fNIRS to examine neural responses to real versus virtual emotional faces, however, several limitations should be acknowledged.
First, since we recruited only from Ontario Tech University's undergraduate student body, our sample of students were limited to that relatively young age range and community, and likely are more educated and technologically savvy than the general population.
Future studies should include a wider range of participants, ideally representing the general population. 
Second, since the virtual faces used in this study were all from one dataset (UIBVFED), they were all of similar realism and stylization.
Future work should explore a wider range of virtual face styles and realism levels to better understand how these factors influence neural processing, and generalize across different virtual face types \citep{schindler_differential_2017}.

\section*{Conclusion}
The present study investigated how the human brain processes emotional expressions on real versus virtual faces, using functional near-infrared spectroscopy (fNIRS) to assess both activation magnitude and functional connectivity. 
By directly comparing hemodynamic responses to real and virtual faces across a range of basic emotions, we aimed to determine whether face realism modulates neural responses, and whether these effects extend beyond local activation to distributed network interactions. 
We found that virtual faces elicited greater activation in the left occipital region, suggesting increased perceptual demands, while real faces were associated with stronger distributed connectivity and better recall performance in a memory task. 
Neutral and Surprise expressions evoked the strongest activation, whereas Fear and Anger engaged broader neural networks. 
These findings highlight the complex interplay between face realism and emotion processing, suggesting that while virtual faces are perceived as faces, they do not recruit the same neural mechanisms as real ones. 
Emotional expressions appear to be decoded in distributed brain networks rather than isolated regions. 
These results offer novel insights into the neural mechanisms of emotional face perception in the context of increasingly prevalent virtual social interactions. 
Future research should examine how varying levels of virtual face realism and stylization influence neural responses and social cognition.